\section{Vatican I: What the Council Actually Promulgated}

The First Vatican Council opened in St Peter's Basilica on 8 December 1869. More than seven hundred fathers participated at some point, overwhelmingly Latin-rite bishops but with Eastern Catholic hierarchs and superiors general also present. No separated Christian observers were invited. The scale was unprecedented, while the procedures remained heavily Roman: preparatory commissions had drafted the texts; conciliar deputies could revise them; the pope controlled the agenda and reserved final approval. The minority was genuine and articulate, but “majority” and “minority” changed by question. Many who opposed defining infallibility as inopportune did not deny papal primacy or the underlying doctrine.

\subsection{A council of two constitutions}

The public result was much smaller than the planned agenda.

\begin{threestable}{Promulgated act}{Doctrinal architecture}{What it accomplished}
\emph{Dei Filius}, 24 April 1870 & God the Creator; Revelation; faith; faith and reason; eighteen canons with anathemas. & Defined natural knowledge of God, supernatural Revelation and faith, the objective content of faith, the Church's credibility, and the harmony and proper distinction of faith and reason; rejected pantheism, materialism, rationalism, fideist errors, and relativizing development.\\
\emph{Pastor aeternus}, 18 July 1870 & Petrine institution; perpetuity in Roman pontiffs; universal primacy of jurisdiction; infallible magisterium. & Defined the Roman primacy as immediate, ordinary, episcopal, and universal while affirming ordinary episcopal authority; then defined the conditions and effect of ex cathedra judgments.\\
Unfinished schemas & Church, episcopate, clergy, missions, religious life, discipline, and other proposed matters. & Debated or prepared but never promulgated; they cannot be cited as acts of Vatican I.\\
\end{threestable}

Calling Vatican I “the infallibility council” therefore hides half of its dogmatic achievement and most of its intended scope. \emph{Dei Filius} supplies the epistemology that the papal definition presupposes: Revelation is given by God, the deposit is objectively received, faith is not a private sentiment, reason is real, and the magisterium serves rather than creates divine truth.

\subsection{The four chapters of \emph{Dei Filius}}

Chapter 1 confesses one true and living God, Creator of all things, distinct from the world. Creation is free, not a necessary emanation of divine substance. Chapter 2 teaches that God can be known with certainty by natural reason from created things, while supernatural Revelation makes divine truths accessible to all with firmness and without admixture of error. This is not a promise that every person easily reaches a complete natural theology; it identifies the capacity and objective knowability of God.

Chapter 3 treats faith as a supernatural virtue by which, moved and assisted by grace, a person believes revealed things because of the authority of God revealing. Motives of credibility—miracles, prophecy, the Church's marks—show that assent is not a blind movement, while grace keeps it from being merely a conclusion of natural evidence. The chapter's statement that those who have received the faith under the Church's magisterium can never have a just cause for changing or doubting it concerns the objective warrants of revealed faith; it does not deny psychological trial or make coercion an act of faith.

Chapter 4 refuses both conflict and confusion. The same God gives reason and reveals mysteries. A real contradiction cannot exist, though an opinion misnamed science or a misunderstood dogma can appear to conflict. Reason serves faith by demonstrating preambles, illuminating analogies, and connecting mysteries; it never comprehends God or converts mystery into autonomous proof. Dogmatic development occurs within the same meaning and judgment, a phrase inherited from Vincent of Lérins. Growth is real; mutation into a contradictory sense is not.

\subsection{From the Church schema to \emph{Pastor aeternus}}

The initial schema on the Church treated many subjects. External events and the organized desire for a definition brought the papal chapters forward. Debate addressed not only whether the pope is infallible but whether a definition was opportune, how the Church's infallibility relates to his, whether consultation should be named, which objects are included, what “define” means, and whether subsequent consent is required.

The conciliar minority included figures such as Karl Joseph von Hefele, Félix Dupanloup, and several Eastern hierarchs. Their reasons varied: historical cases including Honorius I, fear of ecumenical damage, defense of episcopal witness, objection to the timing, and concern that public maximalism would swallow every papal act. Some proposed language requiring the counsel or testimony of the Churches. The deputation's spokesman, Bishop Vincent Gasser, answered these objections in his \emph{relatio} of 11 July. That explanation is not itself the dogmatic definition, but it is privileged evidence of what the deputation and voting fathers were asked to mean.

On 13 July, in the decisive general-congregation vote on the revised schema, 451 fathers voted \emph{placet}, 62 \emph{placet iuxta modum}, and 88 \emph{non placet}. Many opponents then left Rome with permission rather than vote against the pope in the public session. On 18 July, 535 fathers voted: 533 in favor and two against; the two immediately submitted. The overwhelming final number must not be used to erase the recorded objections or pretend the preceding vote was unanimous. Neither should the minority's departure be converted into a veto that the constitution itself denied.

\subsection{The whole teaching of \emph{Pastor aeternus}}

Chapter 1 defines that Christ immediately and directly conferred on Peter a true primacy of jurisdiction over the whole Church. Chapter 2 teaches its perpetuity in the Roman pontiffs. Chapter 3 describes the pope's jurisdiction as full, supreme, ordinary, immediate, and episcopal over pastors and faithful, while expressly stating that this power does not harm the ordinary and immediate episcopal jurisdiction by which bishops govern their assigned flocks. Recourse to the Roman See is final; civil government may not obstruct ecclesial communication.

The fourth chapter first treats the Holy See's doctrinal service in the history of the Church and then defines one specifically conditioned exercise of infallibility. The famous final sentence cannot be detached from the prologue's purpose—unity of faith and communion—or from \emph{Dei Filius}'s account of Revelation. Nor may papal infallibility be used to eclipse the more continuously exercised and institutionally consequential primacy of jurisdiction defined in chapter 3.

\subsection{War, Rome, and an unfinished council}

France declared war on Prussia on 19 July, the day after \emph{Pastor aeternus}. Napoleon III withdrew the troops protecting Rome. Italian forces entered the city on 20 September. Pius IX suspended the council indefinitely by \emph{Postquam Dei munere} on 20 October 1870. It was never reconvened and was formally closed only in the sense that John XXIII later convoked a new council, not a resumed fifth session.

The interruption is not a defect in the validity of the two constitutions. It is an interpretive fact. A “First Dogmatic Constitution on the Church of Christ” was promulgated without the contemplated second constitution on the Church, episcopate, and other members. The papal head was described at length before the episcopal body and whole people could receive equivalent conciliar treatment. Vatican II would deliberately take up that unfinished ecclesiology.

\judgmentbox{Vatican I did not pronounce that every papal statement is infallible, that the pope receives new revelation, that bishops are papal branch managers, or that faith defeats reason. It promulgated two tightly argued dogmatic constitutions. Its historical imbalance lies not in a false definition but in a true papal doctrine left standing without the council's intended full account of the Church.}
